\section{Introduction}
\label{sec:intro}

Matrix-valued states appear in fast-weight architectures
\citep{schlag2021linear} and extensions of continuous chain-of-thought
models \citep{hao2024coconut}. Their rank describes the dimension of the
image of the linear transformation they represent. Does training learn
the rank needed for a task, and does restricting that rank impair
performance?

A recent study added a $d \times d$ matrix latent to a vector-pretrained
GPT-2 reasoner and found flat rank-$k$ truncation curves across four
training regimes and four readouts \citep{larson2026gradient}. The study
identified a limitation: its ProsQA task admits a rank-1 solution.
That result therefore does not establish how training behaves when
exact recovery requires higher rank.

We study models trained from scratch on tasks whose exact linear
solutions require $\rank(Z) \geq K$. A fixed readout and a single matrix
state isolate the role of rank (\S\ref{sec:setup}). We measure the rank
learned during training and the effect of training-time rank caps
(\S\ref{sec:recruit}). We then test repeated application of the learned
operator, analyze its action on the entity subspace
(\S\ref{sec:compose}), and examine recovery at larger matrix dimensions
and longer training budgets (\S\ref{sec:frontier}).
